\SourceLeaf{228}%
\[
\cos(1)L=X,\qquad\cos(2)L=Y,\qquad\cos(3)L=Z
\]
\noindent and we denote the coordinates of the point \(A\) by \(x,y,z\). Let \(x+\mathrm{d}x,y+\mathrm{d}y,z+\mathrm{d}z\) be the coordinates of another point \(A'\) on the curved surface. Let \(\mathrm{d}s\) be its infinitely small distance from \(A\), and let \(\lambda\) be the point on the spherical surface that represents the direction of the element \(AA'\). Then
\[
\mathrm{d}x=\mathrm{d}s\mathord{.}\cos(1)\lambda,\quad\mathrm{d}y=\mathrm{d}s\mathord{.}\cos(2)\lambda,\quad\mathrm{d}z=\mathrm{d}s\mathord{.}\cos(3)\lambda
\]
\noindent and since we must have \(\lambda L=90^0\),
\[
X\cos(1)\lambda+Y\cos(2)\lambda+Z\cos(3)\lambda=0
\]
\noindent Combining these equations, we obtain
\[
X\mathrm{d}x+Y\mathrm{d}y+Z\mathrm{d}z=0
\]
\par\indent There are two general methods for specifying a curved surface. The \textit{first} uses an equation between the coordinates \(x,y,z\), which we assume has been reduced to the form \(W=0\), where \(W\) is a function of the indeterminates \(x,y,z\). Let the total differential of the function \(W\) be
\[
\mathrm{d}W=P\mathrm{d}x+Q\mathrm{d}y+R\mathrm{d}z
\]
\noindent and on the curved surface we have
\[
P\mathrm{d}x+Q\mathrm{d}y+R\mathrm{d}z=0
\]
\noindent and therefore
\[
P\cos(1)\lambda+Q\cos(2)\lambda+R\cos(3)\lambda=0
\]
\noindent Since this equation, like the one we established above, must hold for the directions of all elements \(\mathrm{d}s\) on the curved surface, we can readily see that \(X,Y,Z\) must be proportional to \(P,Q,R\). Consequently, from
\[
XX+YY+ZZ=1
\]
\noindent we obtain either
\begingroup\small
\[
X=\frac{P}{\BareRoot{(PP+QQ+RR)}},\quad Y=\frac{Q}{\BareRoot{(PP+QQ+RR)}},\quad Z=\frac{R}{\BareRoot{(PP+QQ+RR)}}
\]\endgroup

\noindent or
\begingroup\small
\[
X=\frac{-P}{\BareRoot{(PP+QQ+RR)}},\quad Y=\frac{-Q}{\BareRoot{(PP+QQ+RR)}},\quad Z=\frac{-R}{\BareRoot{(PP+QQ+RR)}}
\]\endgroup

